\documentclass[11pt,a4paper,oneside]{book}
\usepackage[T1]{fontenc}
\usepackage[utf8]{inputenc}
\usepackage{lmodern}
\usepackage{microtype}
\usepackage{amsmath,amssymb,amsfonts,mathtools}
\usepackage{bm}
\usepackage{physics}
\usepackage{tensor}
\usepackage{siunitx}
\usepackage{graphicx}
\usepackage{booktabs}
\usepackage{array}
\usepackage{longtable}
\usepackage{multirow}
\usepackage{tikz}
\usetikzlibrary{arrows.meta,positioning,calc,decorations.pathreplacing}
\usepackage{xcolor}
\usepackage{enumitem}
\usepackage{hyperref}
\usepackage{cleveref}
\usepackage[numbers,sort&compress]{natbib}
\usepackage{geometry}
\geometry{a4paper,left=30mm,right=25mm,top=28mm,bottom=30mm}
\usepackage{makeidx}
\makeindex

\newtheorem{principle}{Principle}[chapter]
\newtheorem{remark}{Remark}[chapter]
\newtheorem{definition}{Definition}[chapter]

\newcommand{\R}{\mathbb{R}}
\newcommand{\C}{\mathbb{C}}
\newcommand{\N}{\mathbb{N}}
\newcommand{\T}{\mathsf{T}}
\newcommand{\Loss}{\mathcal{L}}
\newcommand{\D}{\mathrm{d}}
\newcommand{\E}{\mathbb{E}}
\newcommand{\Prob}{\mathbb{P}}
\newcommand{\ML}{\mathrm{ML}}
\newcommand{\AI}{\mathrm{AI}}
\newcommand{\PINN}{\mathrm{PINN}}
\newcommand{\PDE}{\mathrm{PDE}}
\newcommand{\EFT}{\mathrm{EFT}}

\title{
\Huge\bfseries
Artificial Intelligence and Physics\\[0.4cm]
\Large
From Analytical and Numerical Methods\\
to Machine-Assisted Scientific Discovery
}

\author{
}

\date{\today}

% ============================================================
\begin{document}
% ============================================================

\frontmatter

\maketitle

\chapter*{Preface}

Artificial intelligence has entered physics at a remarkable rate. Machine
learning is now routinely used for experimental classification, parameter
estimation, simulation, inverse problems, optimization, and increasingly for
the discovery of mathematical relations. Large language models have added
another layer: they can manipulate scientific literature, mathematical
expressions, computer programs, and verbal descriptions of theories.

It is tempting to interpret these developments as the beginning of a
replacement of theoretical physics by artificial intelligence. Such an
interpretation is premature.

The purpose of this book is to examine a more fundamental question:

\begin{quote}
What happens to the methodology of physics when artificial intelligence becomes
an active participant in the scientific process?
\end{quote}

The answer requires distinguishing several activities that are often conflated.
Physics uses analytical reasoning to expose structure, numerical methods to
obtain quantitative solutions when analytical methods fail, experiments to
confront theories with nature, and statistical methods to infer physical
information from imperfect observations. Machine learning introduces a new
possibility: learning approximations, representations, operators, or even
candidate equations directly from data and simulations.

The central thesis of this book is therefore not

\[
\text{AI replaces physics},
\]

but rather

\[
\boxed{
\text{analytical physics}
+
\text{numerical physics}
+
\text{AI}
+
\text{experiment}
}
\]

as a new integrated methodology.

The distinction between prediction and understanding will be particularly
important. An algorithm may predict a physical observable with extraordinary
accuracy while providing little information about the reason for the
regularity. Conversely, a compact analytical theory may provide enormous
conceptual understanding while being incapable of producing a numerical answer
for a particular complicated system.

AI changes this balance. It may make previously inaccessible calculations
possible, search spaces too large for human exploration, discover empirical
regularities, and propose mathematical structures. But the resulting candidates
still require physical interpretation and independent verification.

The book consequently treats AI neither as an oracle nor as a threat to
physics, but as a new scientific instrument.

\tableofcontents

\listoffigures
\listoftables

% ============================================================
\mainmatter
% ============================================================

\part{The Methodology of Physics}

% ============================================================
\chapter{Physics Before Artificial Intelligence}
\label{chap:before-ai}
% ============================================================

\section{Physics as a hierarchy of descriptions}

Physical science operates simultaneously at several levels.

At the empirical level one has observations,

\begin{equation}
D=\{x_1,x_2,\ldots,x_N\}.
\end{equation}

At the phenomenological level one constructs relations among observables,

\begin{equation}
y=f(x_1,\ldots,x_n).
\end{equation}

At the theoretical level one introduces mathematical structures such as
differential equations, variational principles, Hamiltonians, Lagrangians and
field equations.

At the deepest level one searches for principles such as symmetry, locality,
causality, conservation laws and invariance.

The methodology of physics can therefore be schematically represented as

\begin{equation}
\text{observation}
\longrightarrow
\text{regularity}
\longrightarrow
\text{model}
\longrightarrow
\text{theory}
\longrightarrow
\text{prediction}
\longrightarrow
\text{experiment}.
\end{equation}

Artificial intelligence enters primarily between observation and model, but it
can potentially participate in almost every stage.

\section{Analytical physics}

Analytical methods seek structures that can be manipulated symbolically.

Examples include

\begin{itemize}
\item symmetry arguments,
\item dimensional analysis,
\item conservation laws,
\item variational principles,
\item perturbation theory,
\item asymptotic analysis,
\item exact solutions,
\item canonical transformations,
\item group-theoretical methods,
\item effective field theory.
\end{itemize}

Their principal advantage is not necessarily numerical accuracy. It is
understanding.

For example, the equation

\begin{equation}
\partial_\mu T^{\mu\nu}=0
\end{equation}

contains structural information that a table of numerical values does not.

\section{Numerical physics}

Most realistic physical systems cannot be solved analytically.

One therefore replaces

\begin{equation}
\mathcal{F}[u]=0
\end{equation}

by a finite-dimensional approximation

\begin{equation}
\mathcal{F}_h(\mathbf{u}_h)=0.
\end{equation}

The subscript $h$ represents a discretization scale.

The numerical problem may then be solved by linear algebra, iteration,
optimization, Monte Carlo sampling, spectral methods, finite differences,
finite elements, or other techniques.

Numerical physics therefore supplies quantitative information where analytical
methods become insufficient.

\section{The emergence of machine learning}

Machine learning introduces a different route:

\begin{equation}
\{x_i,y_i\}_{i=1}^N
\longrightarrow
f_\theta,
\end{equation}

where $\theta$ is determined from the data.

The conceptual transition is

\begin{equation}
\boxed{
\text{derive the algorithm from the equations}
}
\end{equation}

to

\begin{equation}
\boxed{
\text{learn the approximation from examples}.
}
\end{equation}

Modern scientific machine learning combines both approaches.

% ============================================================
\chapter{Three Computational Paradigms}
% ============================================================

\section{Analytical, numerical and AI descriptions}

Consider a physical problem

\begin{equation}
\mathcal{F}[u;\lambda]=0,
\end{equation}

where $\lambda$ denotes physical parameters.

There are three basic strategies.

\subsection{Analytical strategy}

Find

\begin{equation}
u=u(x;\lambda)
\end{equation}

in closed or controlled analytical form.

\subsection{Numerical strategy}

Construct

\begin{equation}
u_h(x;\lambda)
\end{equation}

by discretization.

\subsection{Machine-learning strategy}

Construct

\begin{equation}
u_\theta(x;\lambda)
\end{equation}

by optimization over a parameterized function space.

The three approaches should not be considered competitors in a simple sense.

\begin{equation}
\boxed{
u_{\rm analytic}
\quad
u_h
\quad
u_\theta
}
\end{equation}

are different representations of the same physical object.

\section{The hierarchy of approximation}

Every computational description introduces approximation.

Analytical approximation may involve

\begin{equation}
\epsilon_{\rm asymptotic},
\end{equation}

numerical approximation

\begin{equation}
\epsilon_{\rm discretization},
\end{equation}

and machine learning

\begin{equation}
\epsilon_{\rm approximation}
+
\epsilon_{\rm optimization}
+
\epsilon_{\rm generalization}.
\end{equation}

A complete error budget therefore has the schematic form

\begin{equation}
\epsilon_{\rm total}
=
\epsilon_{\rm model}
+
\epsilon_{\rm numerical}
+
\epsilon_{\rm statistical}
+
\epsilon_{\rm ML}.
\end{equation}

This is one reason why AI cannot simply be inserted into physics without
changing the standards of validation.

\section{What AI adds}

AI is particularly useful when

\begin{enumerate}
\item the input dimension is high;
\item the forward model is expensive;
\item the data are abundant;
\item the desired mapping is nonlinear;
\item repeated simulations are required;
\item the inverse problem is difficult;
\item the relevant functional form is unknown.
\end{enumerate}

These properties occur throughout modern physics.

% ============================================================
\chapter{Machine Learning as Function Approximation}
% ============================================================

\section{The basic problem}

Let

\begin{equation}
f:\mathcal{X}\rightarrow\mathcal{Y}
\end{equation}

be an unknown mapping.

Given data

\begin{equation}
D=
\{(x_i,y_i)\}_{i=1}^N,
\end{equation}

we seek

\begin{equation}
f_\theta(x)\simeq f(x).
\end{equation}

The parameters are usually obtained from

\begin{equation}
\theta^\ast
=
\operatorname*{arg\,min}_{\theta}
\Loss(\theta).
\end{equation}

For regression,

\begin{equation}
\Loss_{\rm data}
=
\frac{1}{N}
\sum_{i=1}^N
\left\|
f_\theta(x_i)-y_i
\right\|^2.
\end{equation}

This apparently simple equation is the basis of a large fraction of
machine-learning applications in physics.

\section{Universal approximation versus physical approximation}

Neural networks possess powerful approximation properties. But the statement

\begin{equation}
f_\theta\simeq f
\end{equation}

does not imply that

\begin{equation}
f_\theta
\end{equation}

has the same physical structure as $f$.

A network can approximate a function while violating:

\begin{itemize}
\item symmetry,
\item conservation laws,
\item positivity,
\item causality,
\item dimensional consistency,
\item locality.
\end{itemize}

The problem is therefore not approximation alone.

It is

\begin{equation}
\boxed{
\text{physically admissible approximation}.
}
\end{equation}

% ============================================================
\chapter{Physics-Informed Machine Learning}
% ============================================================

\section{From data fitting to constrained learning}

Suppose

\begin{equation}
\mathcal{L}[u]=0
\end{equation}

is the governing differential equation.

A neural network represents

\begin{equation}
u_\theta(x).
\end{equation}

Define the residual

\begin{equation}
r_\theta(x)
=
\mathcal{L}[u_\theta](x).
\end{equation}

Instead of minimizing only data error, introduce

\begin{equation}
\Loss
=
\lambda_d\Loss_d
+
\lambda_p\Loss_p,
\end{equation}

where

\begin{equation}
\Loss_p
=
\frac{1}{N_c}
\sum_{j=1}^{N_c}
|r_\theta(x_j)|^2.
\end{equation}

The points $x_j$ are collocation points.

This is the essential idea of physics-informed neural networks (PINNs).
Physics-informed ML is broader than PINNs and includes architectural,
variational, operator-learning and other approaches. The literature emphasizes
that physical information can be incorporated through the loss function,
network architecture, invariants and other constraints
\citep{Raissi2019,Karniadakis2021}.

\section{Initial and boundary conditions}

For

\begin{equation}
\partial_tu+F(u)=0,
\end{equation}

with

\begin{equation}
u(x,0)=u_0(x),
\end{equation}

one may write

\begin{equation}
\Loss
=
\lambda_f\Loss_f
+
\lambda_i\Loss_i.
\end{equation}

For boundary conditions,

\begin{equation}
u|_{\partial\Omega}=g,
\end{equation}

one adds

\begin{equation}
\Loss_b
=
\frac{1}{N_b}
\sum_j
|u_\theta(x_j)-g(x_j)|^2.
\end{equation}

Thus

\begin{equation}
\Loss_{\rm total}
=
\Loss_{\rm PDE}
+
\Loss_{\rm IC}
+
\Loss_{\rm BC}
+
\Loss_{\rm data}.
\end{equation}

\section{The conceptual significance}

PINNs illustrate a fundamental transformation:

\begin{equation}
\boxed{
\text{equation as algorithm}
\quad\longrightarrow\quad
\text{equation as constraint}.
}
\end{equation}

The differential equation no longer necessarily determines the numerical
algorithm. It can instead constrain a learned representation.

This is a profound change in methodology.

\section{Limitations of PINNs}

PINNs are not universal replacements for numerical PDE solvers.

Important difficulties include:

\begin{itemize}
\item optimization stiffness;
\item poor treatment of multiscale systems;
\item spectral bias;
\item difficult loss balancing;
\item boundary-layer resolution;
\item high-dimensional domains;
\item long-time integration;
\item chaotic dynamics.
\end{itemize}

Consequently, the appropriate conclusion is not that PINNs have replaced
numerical analysis, but that they provide another class of approximation.

% ============================================================
\chapter{Neural Differential Equations}
% ============================================================

\section{Learning dynamics}

Consider

\begin{equation}
\frac{d\mathbf{x}}{dt}
=
F(\mathbf{x},t).
\end{equation}

A neural differential equation replaces or augments $F$ by

\begin{equation}
\frac{d\mathbf{x}}{dt}
=
F_\theta(\mathbf{x},t).
\end{equation}

The network therefore represents the vector field itself.

This differs conceptually from learning only the trajectory.

\section{Learning unknown physics}

Suppose

\begin{equation}
\frac{d x}{dt}
=
F_{\rm known}(x)
+
F_{\rm unknown}(x).
\end{equation}

One may construct

\begin{equation}
\frac{dx}{dt}
=
F_{\rm known}(x)
+
F_\theta(x).
\end{equation}

This hybrid construction is particularly attractive scientifically because
known physics remains explicit.

It represents

\begin{equation}
\boxed{
\text{known theory}
+
\text{learned correction}.
}
\end{equation}

% ============================================================
\chapter{Neural Operators}
% ============================================================

\section{From functions to functions}

A conventional network approximates

\begin{equation}
f:\R^n\rightarrow\R^m.
\end{equation}

A physical field theory frequently requires

\begin{equation}
\mathcal{G}:u\mapsto v,
\end{equation}

where $u$ and $v$ are functions.

For example,

\begin{equation}
-\nabla^2u=f
\end{equation}

defines an operator

\begin{equation}
\mathcal{G}:f\mapsto u.
\end{equation}

Neural operators seek

\begin{equation}
\mathcal{G}_\theta\simeq\mathcal{G}.
\end{equation}

Recent work emphasizes neural operators precisely as a framework for learning
mappings between function spaces, with applications to PDEs, fluid dynamics,
materials and inverse design \citep{Azizzadenesheli2024}.

\section{Why operator learning matters}

Suppose the same PDE must be solved for many initial conditions

\begin{equation}
u_0^{(1)},u_0^{(2)},\ldots,u_0^{(N)}.
\end{equation}

A conventional solver computes each trajectory separately.

An operator-learning system instead approximates

\begin{equation}
u(t)
=
\mathcal{G}_t[u_0].
\end{equation}

After training,

\begin{equation}
u_\theta(t)
=
\mathcal{G}_{\theta,t}[u_0]
\end{equation}

may be evaluated rapidly.

This is one of the strongest cases in which AI can genuinely change the
computational methodology.

\section{Fourier neural operators}

A Fourier neural operator represents transformations in Fourier space.

Schematically,

\begin{equation}
v_{l+1}
=
\sigma
\left(
W_l v_l
+
\mathcal{F}^{-1}
\left[
K_l(k)\,
\mathcal{F}[v_l]
\right]
\right).
\end{equation}

The kernel acts on spectral modes.

This is particularly natural for translationally structured systems and
periodic domains.

\section{Operator learning versus PDE solving}

It is important to distinguish

\begin{equation}
\text{PDE solver}
\end{equation}

from

\begin{equation}
\text{learned solution operator}.
\end{equation}

The first solves a specified equation.

The second learns a family of solutions generated by a class of equations or
parameters.

Thus the object of computation changes.

% ============================================================
\chapter{Symmetry, Invariance and Equivariance}
% ============================================================

\section{Symmetry in physics}

Suppose a group $G$ acts on the configuration space:

\begin{equation}
x\mapsto g x.
\end{equation}

An invariant observable satisfies

\begin{equation}
f(gx)=f(x).
\end{equation}

An equivariant map satisfies

\begin{equation}
F(gx)=gF(x),
\end{equation}

or more generally

\begin{equation}
F(\rho_X(g)x)
=
\rho_Y(g)F(x).
\end{equation}

\section{Why symmetry should not always be learned}

If rotational invariance is known exactly, asking a network to infer it from
data is inefficient.

Instead, the architecture can impose

\begin{equation}
f(Rx)=f(x).
\end{equation}

The network then learns only the physically independent information.

\section{Symmetry as data compression}

Symmetry reduces the effective space of configurations.

If a group orbit is

\begin{equation}
\mathcal{O}_x
=
\{gx\,|\,g\in G\},
\end{equation}

then all points in the orbit may represent equivalent physical situations.

Learning on the quotient space

\begin{equation}
\mathcal{X}/G
\end{equation}

is therefore conceptually preferable.

\section{Gauge symmetry}

Gauge symmetry is more subtle.

A gauge transformation

\begin{equation}
A_\mu
\rightarrow
A_\mu+\partial_\mu\Lambda
\end{equation}

does not represent a new physical configuration.

A machine-learning representation of gauge fields should therefore distinguish
gauge redundancy from physical degrees of freedom.

This is a particularly important example where naive data-driven learning can
be wasteful or even misleading.

% ============================================================
\chapter{Conservation Laws and Variational Structure}
% ============================================================

\section{Conservation laws}

Suppose

\begin{equation}
\frac{dQ}{dt}=0.
\end{equation}

An unconstrained learned dynamics

\begin{equation}
\dot{x}=F_\theta(x)
\end{equation}

will generally not conserve $Q$ exactly.

One may impose

\begin{equation}
\nabla Q(x)\cdot F_\theta(x)=0.
\end{equation}

The conservation law is then built directly into the learning problem.

\section{Hamiltonian systems}

Hamiltonian mechanics has

\begin{equation}
\dot{q}
=
\frac{\partial H}{\partial p},
\qquad
\dot{p}
=
-\frac{\partial H}{\partial q}.
\end{equation}

Instead of learning the vector field independently, one may learn

\begin{equation}
H_\theta(q,p)
\end{equation}

and define the dynamics by

\begin{equation}
\dot{q}
=
\frac{\partial H_\theta}{\partial p},
\qquad
\dot{p}
=
-\frac{\partial H_\theta}{\partial q}.
\end{equation}

This automatically preserves the Hamiltonian structure.

\section{Lagrangian learning}

Similarly, suppose

\begin{equation}
S[q]
=
\int L(q,\dot q,t)\,dt.
\end{equation}

The Euler--Lagrange equations are

\begin{equation}
\frac{d}{dt}
\frac{\partial L}{\partial\dot q}
-
\frac{\partial L}{\partial q}
=0.
\end{equation}

AI can therefore be used to search over candidate Lagrangians rather than
directly fitting accelerations.

This is scientifically important because the Lagrangian contains much more
structure than a list of trajectories.

\section{Noether structure}

If

\begin{equation}
\delta L=0
\end{equation}

under a continuous transformation, Noether's theorem produces a conserved
quantity.

A machine-learning system that learns the Lagrangian can therefore potentially
recover both

\begin{equation}
\text{dynamics}
\end{equation}

and

\begin{equation}
\text{conservation laws}.
\end{equation}

This is much closer to theoretical physics than direct trajectory fitting.

% ============================================================
\chapter{Inverse Problems}
% ============================================================

\section{Forward versus inverse problems}

The forward problem is

\begin{equation}
u
\longrightarrow
y=\mathcal{F}[u].
\end{equation}

The inverse problem is

\begin{equation}
y
\longrightarrow
u.
\end{equation}

The inverse map may not exist uniquely.

\section{Ill-posedness}

Hadamard's conditions require:

\begin{enumerate}
\item existence,
\item uniqueness,
\item continuous dependence on data.
\end{enumerate}

Inverse physical problems often violate one or more of these conditions.

Machine learning does not remove the ill-posedness.

Instead it supplies a prior through the learned representation.

Thus

\begin{equation}
u_\theta
\end{equation}

should be understood as a regularized solution, not necessarily the unique
physical solution.

\section{Bayesian interpretation}

A principled formulation is

\begin{equation}
p(u|D)
\propto
p(D|u)p(u).
\end{equation}

Machine learning can approximate either

\begin{equation}
p(D|u)
\end{equation}

or

\begin{equation}
p(u|D).
\end{equation}

This connects AI naturally to statistical physics and Bayesian inference.

% ============================================================
\chapter{Surrogate Models and Numerical Simulation}
% ============================================================

\section{The expensive forward model}

Suppose

\begin{equation}
y=F(\lambda)
\end{equation}

requires an expensive numerical simulation.

A surrogate learns

\begin{equation}
F_\theta(\lambda)\simeq F(\lambda).
\end{equation}

The cost is then shifted from repeated simulation to one training phase.

\section{Parameter scans}

This is particularly useful when

\begin{equation}
\lambda\in\R^n
\end{equation}

must be scanned over a large parameter region.

Applications include:

\begin{itemize}
\item cosmological parameter inference;
\item particle-physics phenomenology;
\item materials design;
\item plasma simulations;
\item fluid dynamics;
\item climate modelling.
\end{itemize}

\section{The danger of emulator confidence}

A surrogate may have small error on

\begin{equation}
\lambda\in\mathcal{D}_{\rm train}
\end{equation}

and enormous error outside it.

Therefore one must distinguish

\begin{equation}
\text{accuracy}
\end{equation}

from

\begin{equation}
\text{domain of validity}.
\end{equation}

% ============================================================
\chapter{Symbolic Regression and Equation Discovery}
% ============================================================

\section{From fitting to discovering the functional form}

Ordinary regression assumes a model class.

For example,

\begin{equation}
y=ax+b.
\end{equation}

Symbolic regression instead searches for $f$ itself:

\begin{equation}
y=f(x_1,\ldots,x_n).
\end{equation}

The search space may contain

\begin{equation}
+,\quad -,\quad \times,\quad /,\quad
\sin,\quad \cos,\quad \exp,\quad \log,\ldots
\end{equation}

and possibly derivatives and integrals.

\section{Complexity versus accuracy}

A useful objective is

\begin{equation}
\mathcal{J}(f)
=
\chi^2(f)
+
\lambda C(f),
\end{equation}

where $C(f)$ measures complexity.

This formalizes a version of Occam's principle.

\section{Dimensional consistency}

A candidate expression should satisfy

\begin{equation}
[f]=[y].
\end{equation}

This simple condition eliminates enormous numbers of mathematically possible
but physically meaningless expressions.

Recent work has combined neural symbolic models with Monte Carlo tree search,
genetic programming and language models to improve symbolic discovery of
physical formulas \citep{Ying2025}.

\section{From formula to law}

Finding

\begin{equation}
y=f(x)
\end{equation}

is not yet discovering a physical law.

A physical law should ideally possess:

\begin{enumerate}
\item reproducibility;
\item dimensional consistency;
\item symmetry;
\item appropriate limiting behavior;
\item predictive power;
\item robustness under new experiments;
\item connection to a wider theoretical framework.
\end{enumerate}

Thus

\begin{equation}
\boxed{
\text{equation discovery}
\neq
\text{theory discovery}.
}
\end{equation}

% ============================================================
\chapter{AI and Quantum Mechanics}
% ============================================================

\section{Wavefunction representations}

A quantum state may be written

\begin{equation}
|\psi\rangle
=
\sum_x\psi(x)|x\rangle.
\end{equation}

A neural representation gives

\begin{equation}
\psi(x)\simeq\psi_\theta(x).
\end{equation}

This idea underlies neural-network quantum states and variational approaches
\citep{Carleo2017,Carleo2019}.

\section{Variational principle}

For a Hamiltonian $H$,

\begin{equation}
E[\psi]
=
\frac{
\langle\psi|H|\psi\rangle
}{
\langle\psi|\psi\rangle
}.
\end{equation}

The ground-state energy satisfies

\begin{equation}
E_0
=
\min_\psi E[\psi].
\end{equation}

A neural state therefore provides a variational family

\begin{equation}
\psi_\theta
\end{equation}

and one minimizes

\begin{equation}
E(\theta)
=
\frac{
\langle\psi_\theta|H|\psi_\theta\rangle
}{
\langle\psi_\theta|\psi_\theta\rangle
}.
\end{equation}

The important point is that the quantum variational principle supplies the
physics, while AI supplies the flexible representation.

\section{Quantum state tomography}

Given measurements

\begin{equation}
D=\{M_i,p_i\},
\end{equation}

one may reconstruct

\begin{equation}
\rho.
\end{equation}

Machine learning can represent the density matrix or probability distribution
and infer it from incomplete measurements.

Again, the inverse problem is fundamentally constrained by quantum mechanics.

% ============================================================
\chapter{AI and Quantum Field Theory}
% ============================================================

\section{Lattice field theory}

Lattice calculations involve integrals of the form

\begin{equation}
Z
=
\int\mathcal{D}\phi\,
e^{-S[\phi]}.
\end{equation}

Sampling configurations according to

\begin{equation}
P[\phi]
\propto
e^{-S[\phi]}
\end{equation}

can be computationally expensive.

Machine learning may accelerate sampling or construct approximate generative
models.

\section{Sign problems}

If

\begin{equation}
e^{-S[\phi]}
\end{equation}

is not positive definite, conventional Monte Carlo methods become difficult.

AI-based generative methods may provide alternative representations, although
the fundamental difficulty is not automatically eliminated.

\section{Effective field theory}

Consider

\begin{equation}
\mathcal{L}_{\rm EFT}
=
\mathcal{L}_0
+
\sum_i
\frac{c_i}{\Lambda^{d_i-4}}
\mathcal{O}_i.
\end{equation}

AI can assist in:

\begin{itemize}
\item operator selection;
\item parameter inference;
\item anomaly detection;
\item global scans;
\item identifying experimental signatures.
\end{itemize}

But the operator basis must respect physical principles.

% ============================================================
\chapter{AI in Relativity and Gravitation}
% ============================================================

\section{Einstein's equations}

General relativity is governed by

\begin{equation}
G_{\mu\nu}
+
\Lambda g_{\mu\nu}
=
8\pi G T_{\mu\nu}.
\end{equation}

The nonlinear nature of these equations makes numerical relativity essential
for many problems.

AI may be used for:

\begin{itemize}
\item surrogate waveform generation;
\item parameter estimation;
\item gravitational-wave detection;
\item numerical relativity emulation;
\item solving inverse problems;
\item classification of spacetime structures.
\end{itemize}

\section{The danger of violating covariance}

A learned model must respect the geometric nature of GR.

Tensorial equations must transform appropriately under coordinate changes.

A naive neural network trained in one coordinate system has no automatic reason
to respect general covariance.

This is a particularly clear example of why physical structure must be built
into AI systems.

% ============================================================
\chapter{AI in Condensed Matter and Statistical Physics}
% ============================================================

\section{Phase identification}

Suppose the order parameter is

\begin{equation}
M
=
\frac{1}{N}
\sum_i s_i.
\end{equation}

Different phases may correspond to different statistical distributions of
configurations.

Machine learning can classify phases without explicitly providing the order
parameter.

This raises an interesting possibility:

\begin{equation}
\text{AI classification}
\longrightarrow
\text{candidate order parameter}.
\end{equation}

\section{Critical phenomena}

Near a critical point,

\begin{equation}
\xi\sim |T-T_c|^{-\nu}.
\end{equation}

Machine learning may detect critical behavior from microscopic configurations.

However, identifying $T_c$ statistically is not equivalent to deriving the
critical exponent $\nu$.

Again,

\begin{equation}
\text{pattern recognition}
\neq
\text{theoretical explanation}.
\end{equation}

% ============================================================
\chapter{AI in Cosmology and Astrophysics}
% ============================================================

\section{Parameter inference}

Cosmological observations may be represented by

\begin{equation}
D
=
\{D_1,\ldots,D_N\}.
\end{equation}

The goal is to infer

\begin{equation}
p(\theta|D),
\end{equation}

where $\theta$ contains cosmological parameters.

AI can approximate expensive likelihoods or simulations.

\section{Dark matter}

AI can search for signatures of dark matter in:

\begin{itemize}
\item galaxy rotation curves;
\item gravitational lensing;
\item large-scale structure;
\item particle-detector data.
\end{itemize}

But anomaly detection cannot by itself identify the physical origin of an
anomaly.

% ============================================================
\chapter{AI in Experimental Physics}
% ============================================================

\section{Classification}

Suppose measurements $x$ must be classified into

\begin{equation}
C\in\{1,\ldots,K\}.
\end{equation}

A classifier approximates

\begin{equation}
P(C|x).
\end{equation}

This is one of the most mature applications of AI.

\section{Anomaly detection}

Instead of specifying a signal model, one may learn

\begin{equation}
p(x|\text{known physics})
\end{equation}

and identify events with low likelihood.

This is potentially powerful for new physics.

But

\begin{equation}
\text{anomaly}
\neq
\text{new physics}.
\end{equation}

Detector artifacts, calibration errors and background modelling errors can
produce anomalies.

\section{Experimental design}

Suppose experiment parameters are $\lambda$.

One can optimize

\begin{equation}
\lambda^\ast
=
\operatorname*{arg\,max}_{\lambda}
I(\theta;D|\lambda),
\end{equation}

where $I$ is an information criterion.

AI can therefore participate not merely in analyzing experiments but in
choosing them.

% ============================================================
\chapter{Large Language Models and Theoretical Physics}
% ============================================================

\section{A different type of AI}

Large language models are not ordinary numerical regressors.

They operate on sequences of tokens and learn statistical representations of
language, mathematics, code and other symbolic objects.

Their potential roles include:

\begin{itemize}
\item literature analysis;
\item symbolic manipulation;
\item code generation;
\item theorem exploration;
\item explanation;
\item hypothesis generation;
\item translation between mathematical representations.
\end{itemize}

\section{Mathematical hallucination}

A language model may produce

\begin{equation}
\text{plausible-looking mathematics}
\end{equation}

that is incorrect.

Therefore every important derivation should be checked by:

\begin{enumerate}
\item symbolic algebra;
\item dimensional analysis;
\item limiting cases;
\item numerical substitution;
\item independent derivation.
\end{enumerate}

The relevant distinction is

\begin{equation}
\boxed{
\text{syntactic plausibility}
\neq
\text{mathematical truth}.
}
\end{equation}

\section{LLMs as scientific assistants}

The most realistic near-term role is collaborative.

The physicist formulates the problem and evaluates the result.

The AI assists with exploration and computation.

Thus

\begin{equation}
\text{physicist}
\leftrightarrow
\text{AI}
\leftrightarrow
\text{computer algebra}
\leftrightarrow
\text{numerical solver}.
\end{equation}

% ============================================================
\chapter{Prediction versus Understanding}
% ============================================================

\section{The central epistemological problem}

Consider an AI model

\begin{equation}
y=f_\theta(x)
\end{equation}

that predicts all measured values accurately.

Does it understand physics?

There are at least three levels.

\subsection{Predictive understanding}

The system predicts observations.

\subsection{Structural understanding}

The system identifies mathematical relations and symmetries.

\subsection{Conceptual understanding}

The system identifies mechanisms and principles explaining why the relations
hold.

These are not equivalent.

\section{Example: empirical versus theoretical laws}

Suppose AI discovers

\begin{equation}
F(r)\propto r^{-2}.
\end{equation}

This is a remarkable empirical relation.

But Newtonian gravity contains more:

\begin{equation}
\mathbf{F}
=
-G
\frac{m_1m_2}{r^3}
\mathbf{r}.
\end{equation}

The relation belongs to a dynamical framework.

General relativity then replaces the force description by

\begin{equation}
G_{\mu\nu}
=
8\pi G T_{\mu\nu}.
\end{equation}

Thus successive theories do not merely improve numerical fits. They reorganize
the conceptual structure.

AI may discover the relation without discovering the conceptual reorganization.

% ============================================================
\chapter{Interpretability}
% ============================================================

\section{Why interpretability matters more in physics}

In many engineering applications the question is

\begin{equation}
\text{Does the model work?}
\end{equation}

In fundamental physics one additionally asks

\begin{equation}
\text{Why does it work?}
\end{equation}

Interpretability is therefore not simply a user-interface issue.

It is connected to scientific understanding.

Recent reviews explicitly argue that interpretable ML is central if machine
learning is to contribute to understandable scientific discovery
\citep{Wetzel2025}.

\section{Interpretable representations}

Possible strategies include

\begin{itemize}
\item symbolic regression;
\item sparse models;
\item symmetry-based representations;
\item physically meaningful latent variables;
\item causal models;
\item dimensional reduction.
\end{itemize}

\section{The ideal learned theory}

The ideal result would be

\begin{equation}
D
\longrightarrow
\boxed{
\text{compact mathematical structure}
}
\end{equation}

rather than

\begin{equation}
D
\longrightarrow
\boxed{
\text{opaque numerical predictor}.
}
\end{equation}

% ============================================================
\chapter{Uncertainty and Generalization}
% ============================================================

\section{Prediction is incomplete without uncertainty}

A scientific prediction is usually of the form

\begin{equation}
y=y_0\pm\sigma.
\end{equation}

AI should therefore ideally provide

\begin{equation}
p(y|x,D).
\end{equation}

\section{Sources of uncertainty}

We distinguish:

\begin{enumerate}
\item measurement uncertainty;
\item statistical uncertainty;
\item parameter uncertainty;
\item model uncertainty;
\item numerical uncertainty;
\item extrapolation uncertainty.
\end{enumerate}

\section{Distribution shift}

Training occurs under

\begin{equation}
x\sim p_{\rm train}(x),
\end{equation}

while deployment may involve

\begin{equation}
x\sim p_{\rm test}(x).
\end{equation}

If

\begin{equation}
p_{\rm train}\neq p_{\rm test},
\end{equation}

the learned approximation may fail.

This is particularly dangerous in physics because new discoveries often
involve regimes not present in the training data.

% ============================================================
\chapter{Causality}
% ============================================================

\section{Correlation is not dynamics}

Machine learning excels at estimating

\begin{equation}
P(Y|X).
\end{equation}

Physics is frequently concerned with

\begin{equation}
X(t_0)\longrightarrow X(t).
\end{equation}

The second statement contains causal and dynamical information.

\section{Physical intervention}

A causal model asks what happens if a variable is changed.

This is represented schematically by

\begin{equation}
P(Y|\operatorname{do}(X=x)).
\end{equation}

A purely observational neural network does not automatically recover this
quantity.

Thus causal structure should be incorporated whenever physical interventions
are relevant.

% ============================================================
\chapter{AI and Effective Theories}
% ============================================================

\section{Theory space}

Let $\mathcal{T}$ denote the space of candidate theories.

AI can perform

\begin{equation}
\mathcal{T}
\longrightarrow
\mathcal{T}_{\rm candidate}.
\end{equation}

Physics then imposes

\begin{equation}
\mathcal{T}_{\rm candidate}
\longrightarrow
\mathcal{T}_{\rm consistent}.
\end{equation}

Finally experiment gives

\begin{equation}
\mathcal{T}_{\rm consistent}
\longrightarrow
\mathcal{T}_{\rm viable}.
\end{equation}

\section{Effective field theory as an example}

Consider

\begin{equation}
\mathcal{L}
=
\mathcal{L}_{\rm SM}
+
\sum_i
\frac{c_i}{\Lambda^{d_i-4}}
\mathcal{O}_i.
\end{equation}

AI can explore the coefficient space

\begin{equation}
\{c_i\},
\end{equation}

but the operator basis must respect gauge symmetry and other theoretical
requirements.

This illustrates the complementary roles:

\begin{equation}
\text{AI: search}
\qquad
\text{physics: admissibility}.
\end{equation}

% ============================================================
\chapter{Extrapolation and the Limits of Learning}
% ============================================================

\section{Interpolation}

If

\begin{equation}
x\in\mathcal{D}_{\rm train},
\end{equation}

a model may interpolate successfully.

\section{Extrapolation}

If

\begin{equation}
x\notin\mathcal{D}_{\rm train},
\end{equation}

the model must extrapolate.

Physics is fundamentally interested in extrapolation:

\begin{itemize}
\item high energies;
\item short distances;
\item long times;
\item strong fields;
\item low temperatures;
\item early cosmological epochs.
\end{itemize}

A neural network has no intrinsic reason to extrapolate according to a physical
law.

\section{Physics as an extrapolation prior}

Physical constraints can supply the missing information:

\begin{equation}
\text{data}
+
\text{symmetry}
+
\text{conservation}
+
\text{asymptotics}
\longrightarrow
\text{better extrapolation}.
\end{equation}

This may ultimately be one of the strongest arguments for physics-informed AI.

% ============================================================
\chapter{Verification and Falsifiability}
% ============================================================

\section{The verification problem}

AI can generate candidate hypotheses much faster than humans can verify them.

The scientific bottleneck therefore shifts from

\begin{equation}
\text{hypothesis generation}
\end{equation}

toward

\begin{equation}
\text{hypothesis verification}.
\end{equation}

Recent discussions of AI-driven scientific discovery emphasize precisely this
problem: an abundance of machine-generated hypotheses is useful only if
reliable verification mechanisms scale accordingly.

\section{A verification hierarchy}

Every AI-generated physical result should be subjected, where applicable, to:

\begin{enumerate}
\item dimensional analysis;
\item symmetry tests;
\item conservation laws;
\item limiting cases;
\item analytical benchmarks;
\item numerical benchmarks;
\item statistical validation;
\item independent implementations;
\item experimental tests.
\end{enumerate}

\section{Falsifiability}

A scientific statement must allow the possibility

\begin{equation}
\text{prediction}
\quad\longrightarrow\quad
\text{failure}.
\end{equation}

A flexible AI model that can fit everything may therefore become scientifically
useless.

The goal is not maximal flexibility.

It is

\begin{equation}
\boxed{
\text{maximal explanatory power subject to physical constraints}.
}
\end{equation}

% ============================================================
\chapter{AI and the Discovery of Physical Laws}
% ============================================================

\section{The traditional path}

Historically one may represent discovery as

\begin{equation}
\text{experiment}
\rightarrow
\text{pattern}
\rightarrow
\text{concept}
\rightarrow
\text{equation}.
\end{equation}

AI changes the middle of this process:

\begin{equation}
\text{experiment}
\rightarrow
\text{AI search}
\rightarrow
\text{candidate relations}
\rightarrow
\text{physical interpretation}.
\end{equation}

\section{What would constitute genuine discovery?}

A machine-generated expression

\begin{equation}
F(x)=0
\end{equation}

should be regarded as a candidate law only if it:

\begin{enumerate}
\item survives independent data;
\item has a clear domain of validity;
\item respects exact constraints;
\item has predictive consequences;
\item connects to known physics;
\item survives independent theoretical scrutiny.
\end{enumerate}

\section{Discovery versus rediscovery}

AI may rediscover known laws.

This is scientifically useful because it tests the discovery methodology.

But genuine discovery requires finding structure that was not encoded explicitly
in the training assumptions.

This distinction will become increasingly important as AI systems are trained
on the existing scientific literature.

% ============================================================
\chapter{Autonomous Scientific Discovery}
% ============================================================

\section{From assistant to agent}

The emerging concept of agentic science extends AI beyond isolated tasks.

A possible cycle is

\begin{equation}
\boxed{
\begin{array}{c}
\text{hypothesis generation}\\
\downarrow\\
\text{simulation}\\
\downarrow\\
\text{experiment design}\\
\downarrow\\
\text{measurement}\\
\downarrow\\
\text{analysis}\\
\downarrow\\
\text{hypothesis revision}
\end{array}}
\end{equation}

Recent surveys describe this transition from AI as a computational tool toward
systems capable of combining hypothesis generation, experiment design,
execution and iterative refinement \citep{Wei2025}.

\section{The autonomous laboratory}

A fully integrated system would contain

\begin{equation}
\text{AI}
+
\text{simulation}
+
\text{robotic experiment}
+
\text{statistical inference}.
\end{equation}

Such a system could potentially perform millions of computational experiments
and select a small number for physical verification.

\section{The remaining human role}

The deepest question is not whether the process can be automated.

It is whether the scientific objective itself can be automated.

A system may optimize

\begin{equation}
\mathcal{J}
\end{equation}

without knowing whether $\mathcal{J}$ represents an interesting scientific
question.

Thus human judgment may remain essential at the level of

\begin{equation}
\boxed{
\text{What is worth knowing?}
}
\end{equation}

% ============================================================
\part{AI Across Physics}

% ============================================================
\chapter{Fluid Dynamics}
% ============================================================

\section{Navier--Stokes equations}

The incompressible Navier--Stokes equations are

\begin{equation}
\partial_t\mathbf{u}
+
(\mathbf{u}\cdot\nabla)\mathbf{u}
=
-\frac{1}{\rho}\nabla p
+
\nu\nabla^2\mathbf{u}
+
\mathbf{f},
\end{equation}

with

\begin{equation}
\nabla\cdot\mathbf{u}=0.
\end{equation}

AI can learn:

\begin{itemize}
\item unresolved turbulence;
\item reduced-order dynamics;
\item closure relations;
\item flow reconstruction;
\item control strategies.
\end{itemize}

\section{Hybrid turbulence models}

One may write

\begin{equation}
F_{\rm total}
=
F_{\rm known}
+
F_\theta.
\end{equation}

This provides a principled division between established physics and learned
corrections.

% ============================================================
\chapter{Plasma Physics}
% ============================================================

Plasma physics involves nonlinear kinetic and fluid equations over enormous
ranges of scales.

AI can assist with:

\begin{itemize}
\item turbulence;
\item magnetic confinement;
\item equilibrium reconstruction;
\item instability detection;
\item surrogate simulations;
\item experimental control.
\end{itemize}

The principal challenge is respecting conservation laws and multiscale
structure.

% ============================================================
\chapter{Materials Physics}
% ============================================================

Suppose a material has microscopic structure $X$ and macroscopic property $P$:

\begin{equation}
P=F(X).
\end{equation}

AI can approximate

\begin{equation}
P\simeq F_\theta(X)
\end{equation}

and invert the problem:

\begin{equation}
X\simeq F_\theta^{-1}(P).
\end{equation}

The second relation constitutes inverse materials design.

% ============================================================
\chapter{Climate and Geophysical Physics}
% ============================================================

Climate models contain coupled PDEs with enormous computational costs.

Machine learning can construct

\begin{equation}
F_{\rm climate}
\simeq
F_{\rm resolved}
+
F_{\rm ML},
\end{equation}

where unresolved processes are represented by learned closures.

Recent reviews emphasize the growing role of ML in climate physics while also
highlighting the importance of physical consistency and generalization
\citep{Bracco2025}.

% ============================================================
\part{Foundations and Philosophy}

% ============================================================
\chapter{Can AI Replace Theoretical Physics?}
% ============================================================

The answer depends on what is meant by ``theoretical physics.''

If theory means numerical prediction, then AI can replace parts of traditional
computational methodology.

If theory means compact mathematical description, the answer is less clear.

If theory means conceptual understanding, the answer is presently unknown.

The distinction can be represented as

\begin{equation}
\boxed{
\begin{array}{rcl}
\text{prediction}
&\neq&
\text{model}\\
&\neq&
\text{theory}\\
&\neq&
\text{understanding}.
\end{array}}
\end{equation}

% ============================================================
\chapter{AI as a New Scientific Instrument}
% ============================================================

A telescope does not replace astronomy.

A particle detector does not replace particle physics.

A computer does not replace mathematics.

Likewise,

\begin{equation}
\boxed{
\text{AI need not replace physics}.
}
\end{equation}

It may instead enlarge the observational and computational reach of the
physicist.

This analogy is useful because every scientific instrument changes the
questions that can be asked.

The telescope made questions about distant galaxies meaningful.

The computer made questions about nonlinear many-body dynamics computationally
accessible.

AI may make questions about enormous spaces of models computationally
accessible.

% ============================================================
\chapter{The Human--AI Division of Labor}
% ============================================================

The most plausible future division is:

\begin{center}
\begin{tabular}{p{0.28\textwidth}p{0.62\textwidth}}
\toprule
Human physicist & AI system\\
\midrule
Defines the physical question & Explores large hypothesis spaces\\
Identifies principles & Performs high-dimensional search\\
Chooses relevant observables & Detects statistical structure\\
Interprets concepts & Optimizes representations\\
Judges scientific significance & Generates candidate models\\
Designs validation & Automates repetitive calculations\\
\bottomrule
\end{tabular}
\end{center}

This is not a fixed division. It will evolve.

But the distinction between

\begin{equation}
\text{search}
\end{equation}

and

\begin{equation}
\text{scientific judgment}
\end{equation}

is likely to remain important.

% ============================================================
\chapter{A Unified Methodology}
% ============================================================

We can now formulate a unified scientific workflow.

\begin{equation}
\boxed{
\begin{array}{c}
\text{Physical question}\\
\downarrow\\
\text{Analytical formulation}\\
\downarrow\\
\text{Known constraints}\\
\downarrow\\
\text{Numerical model}\\
\downarrow\\
\text{AI augmentation}\\
\downarrow\\
\text{Prediction}\\
\downarrow\\
\text{Experimental test}\\
\downarrow\\
\text{Theory revision}
\end{array}}
\end{equation}

The key is that AI is embedded in the scientific loop.

It does not stand outside physics.

% ============================================================
\chapter{A Taxonomy of AI in Physics}
% ============================================================

\begin{longtable}{p{0.12\textwidth}p{0.27\textwidth}p{0.48\textwidth}}
\toprule
Level & AI function & Scientific role\\
\midrule
I & Classification & Identify events or phases\\
II & Regression & Estimate observables or parameters\\
III & Inference & Reconstruct hidden variables\\
IV & Surrogate modelling & Replace expensive calculations\\
V & Operator learning & Learn families of physical solutions\\
VI & Anomaly detection & Find deviations from known physics\\
VII & Symbolic regression & Discover candidate equations\\
VIII & Theory exploration & Search candidate theoretical structures\\
IX & Experimental design & Select informative measurements\\
X & Autonomous discovery & Close the entire scientific loop\\
\bottomrule
\end{longtable}

% ============================================================
\chapter{A Critical Assessment}
% ============================================================

\section{What AI does exceptionally well}

AI is particularly strong at:

\begin{enumerate}
\item high-dimensional pattern recognition;
\item interpolation;
\item nonlinear regression;
\item optimization;
\item classification;
\item representation learning;
\item approximate simulation;
\item large-scale search.
\end{enumerate}

\section{What AI does poorly}

Current systems remain vulnerable to:

\begin{enumerate}
\item extrapolation;
\item systematic uncertainty;
\item causal reasoning;
\item exact mathematical consistency;
\item physical interpretability;
\item rare-event reliability;
\item distribution shift.
\end{enumerate}

\section{The fundamental limitation}

The most important limitation is not computational.

It is epistemological.

A model may answer

\begin{equation}
\text{``What happens?''}
\end{equation}

without answering

\begin{equation}
\text{``Why?''}
\end{equation}

Physics has historically regarded the second question as central.

% ============================================================
\chapter{Future Directions}
% ============================================================

Several developments appear especially important.

\section{Theory-aware architectures}

Networks should increasingly encode:

\begin{equation}
\text{symmetry}
+
\text{conservation}
+
\text{geometry}
+
\text{causality}.
\end{equation}

\section{Hybrid symbolic-neural systems}

The combination

\begin{equation}
\text{neural search}
+
\text{symbolic representation}
\end{equation}

is likely to become increasingly important.

\section{AI for mathematical physics}

AI may assist in:

\begin{itemize}
\item conjecture generation;
\item symbolic manipulation;
\item asymptotic analysis;
\item solving differential equations;
\item identifying integrability;
\item exploring representation theory.
\end{itemize}

\section{Automated verification}

As hypothesis generation accelerates, verification must become equally
automated.

The future scientific architecture should therefore include

\begin{equation}
\text{generation}
\longleftrightarrow
\text{verification}.
\end{equation}

% ============================================================
\chapter{Conclusions}
% ============================================================

Artificial intelligence is transforming the methodology of physics.

Its importance does not arise merely because neural networks can approximate
complicated functions. The deeper change is that a physical problem can now be
represented simultaneously in several computational languages:

\begin{equation}
\boxed{
\text{equation}
\quad
\text{numerical algorithm}
\quad
\text{data}
\quad
\text{learned representation}.
}
\end{equation}

The traditional distinction between analytical and numerical physics therefore
acquires a third component.

Analytical methods provide structure.

Numerical methods provide controlled quantitative computation.

AI provides flexible high-dimensional approximation and search.

Experiment provides the final confrontation with nature.

The future methodology is consequently

\begin{equation}
\boxed{
\text{analytical}
+
\text{numerical}
+
\text{AI}
+
\text{experimental}
}.
\end{equation}

The most important conceptual distinction is between prediction and
understanding.

A neural network can reproduce an observable.

A symbolic-regression system can discover a formula.

A neural operator can learn a PDE solution map.

An AI agent can propose an experiment.

None of these achievements, by itself, guarantees a physical theory.

A theory must connect observations through principles, structures and
concepts. It must survive independent tests and make predictions outside the
data from which it was constructed.

Consequently, the proper goal of AI in physics should not be

\begin{equation}
\text{maximum predictive accuracy}
\end{equation}

but

\begin{equation}
\boxed{
\text{maximum scientifically validated explanatory power}.
}
\end{equation}

This suggests a hierarchy:

\begin{equation}
\begin{array}{c}
\text{data}\\
\downarrow\\
\text{prediction}\\
\downarrow\\
\text{pattern}\\
\downarrow\\
\text{mathematical relation}\\
\downarrow\\
\text{physical model}\\
\downarrow\\
\text{theory}\\
\downarrow\\
\text{understanding}.
\end{array}
\end{equation}

AI is already powerful in the lower and middle parts of this hierarchy.

Whether it can reliably participate in the upper part remains an open scientific
question.

The most promising answer is neither optimistic replacement nor pessimistic
rejection.

It is collaboration.

\begin{equation}
\boxed{
\text{Human conceptual physics}
\;\longleftrightarrow\;
\text{AI computational intelligence}
}
\end{equation}

may become one of the defining methodological combinations of twenty-first
century physics.

% ============================================================
\appendix
% ============================================================

\part*{Appendices}
\addcontentsline{toc}{part}{Appendices}

% ============================================================
\chapter{Basic Machine-Learning Mathematics}
% ============================================================

\section{Optimization}

Given

\begin{equation}
\Loss(\theta),
\end{equation}

gradient descent gives

\begin{equation}
\theta_{n+1}
=
\theta_n
-
\eta\nabla_\theta\Loss(\theta_n).
\end{equation}

\section{Gradient-based learning}

For a composite function

\begin{equation}
y=f(g(x)),
\end{equation}

automatic differentiation applies

\begin{equation}
\frac{dy}{dx}
=
\frac{\partial f}{\partial g}
\frac{\partial g}{\partial x}.
\end{equation}

This is central to physics-informed learning because PDE residuals require
derivatives of neural-network outputs.

% ============================================================
\chapter{Dimensional Analysis and AI}
% ============================================================

Suppose there are $n$ dimensional variables and $r$ independent dimensions.

Buckingham's theorem implies the existence of

\begin{equation}
n-r
\end{equation}

independent dimensionless combinations.

Thus a physical relation can be reduced to

\begin{equation}
\Pi_1
=
F(\Pi_2,\ldots,\Pi_{n-r}).
\end{equation}

This provides an immediate way of reducing the dimensionality of AI problems.

% ============================================================
\chapter{A Recommended Workflow for AI-Assisted Physics}
% ============================================================

A robust workflow is:

\begin{enumerate}
\item formulate the physical problem analytically;
\item identify exact constraints;
\item establish known limiting cases;
\item construct a conventional numerical benchmark;
\item define the AI problem;
\item train the model;
\item quantify uncertainty;
\item test interpolation;
\item test extrapolation;
\item compare with analytical limits;
\item compare with numerical solutions;
\item perform independent validation;
\item only then interpret the result physically.
\end{enumerate}

This workflow prevents the AI model from becoming the source of both the
hypothesis and its own validation.

% ============================================================
\chapter{Glossary}
% ============================================================

\begin{description}

\item[Artificial intelligence:]
A broad class of computational methods designed to perform tasks associated
with intelligent behavior.

\item[Machine learning:]
Methods in which model parameters or representations are inferred from data.

\item[Physics-informed machine learning:]
Machine learning in which physical equations, symmetries or constraints are
incorporated explicitly.

\item[PINN:]
Physics-informed neural network.

\item[Neural operator:]
A learned approximation to an operator mapping functions to functions.

\item[Surrogate model:]
A computationally inexpensive approximation to an expensive physical model.

\item[Symbolic regression:]
Search for mathematical expressions directly from data.

\item[Equivariance:]
Transformation of the output consistently with transformation of the input.

\item[Interpretability:]
The ability to understand the internal or functional basis of a model's
prediction.

\item[Distribution shift:]
A difference between the statistical distribution of training and deployment
data.

\item[Autonomous scientific discovery:]
An AI-driven workflow combining hypothesis generation, simulation,
experimentation, analysis and iterative refinement.

\end{description}

% ============================================================
\bibliographystyle{unsrtnat}

\begin{thebibliography}{99}

\bibitem[Carleo and Troyer(2017)]{Carleo2017}
G.~Carleo and M.~Troyer,
``Solving the quantum many-body problem with artificial neural networks,''
Science \textbf{355}, 602--606 (2017).

\bibitem[Carleo et al.(2019)]{Carleo2019}
G.~Carleo, I.~Cirac, K.~Cranmer, L.~Daudet, M.~Schuld,
N.~Tishby, L.~Vogt-Maranto and L.~Zdeborov\'a,
``Machine learning and the physical sciences,''
Rev.\ Mod.\ Phys.\ \textbf{91}, 045002 (2019).

\bibitem[Raissi et al.(2019)]{Raissi2019}
M.~Raissi, P.~Perdikaris and G.~E.~Karniadakis,
``Physics-informed neural networks: A deep learning framework for solving
forward and inverse problems involving nonlinear partial differential
equations,''
J.\ Comput.\ Phys.\ \textbf{378}, 686--707 (2019).

\bibitem[Karniadakis et al.(2021)]{Karniadakis2021}
G.~E.~Karniadakis, I.~G.~Kevrekidis, L.~Lu, P.~Perdikaris,
S.~Wang and L.~Yang,
``Physics-informed machine learning,''
Nature Rev.\ Phys.\ \textbf{3}, 422--440 (2021).

\bibitem[Brunton et al.(2020)]{Brunton2020}
S.~L.~Brunton, B.~R.~Noack and P.~Koumoutsakos,
``Machine learning for fluid mechanics,''
Annu.\ Rev.\ Fluid Mech.\ \textbf{52}, 477--508 (2020).

\bibitem[Karagiorgi et al.(2022)]{Karagiorgi2022}
G.~Karagiorgi, G.~Kasieczka, S.~Kravitz, B.~Nachman and D.~Shih,
``Machine learning in the search for new fundamental physics,''
Nature Rev.\ Phys.\ \textbf{4}, 399--412 (2022).

\bibitem[Azizzadenesheli et al.(2024)]{Azizzadenesheli2024}
K.~Azizzadenesheli, N.~Kovachki, Z.~Li, M.~Liu-Schiaffini,
J.~Kossaifi and A.~Anandkumar,
``Neural operators for accelerating scientific simulations and design,''
Nature Rev.\ Phys.\ \textbf{6}, 320--328 (2024).

\bibitem[Bracco et al.(2025)]{Bracco2025}
A.~Bracco, J.~Brajard, H.~A.~Dijkstra, P.~Hassanzadeh,
C.~Lessig and C.~Monteleoni,
``Machine learning for the physics of climate,''
Nature Rev.\ Phys.\ \textbf{7}, 6--20 (2025).

\bibitem[Wetzel et al.(2025)]{Wetzel2025}
S.~J.~Wetzel, S.~Ha, R.~Iten, M.~Klopotek and Z.~Liu,
``Interpretable Machine Learning in Physics: A Review,''
arXiv:2503.23616 (2025).

\bibitem[Ying et al.(2025)]{Ying2025}
J.~Ying, H.~Lin, C.~Yue et al.,
``A neural symbolic model for space physics,''
Nature Machine Intelligence \textbf{7}, 1726--1741 (2025).

\bibitem[Wei et al.(2025)]{Wei2025}
J.~Wei et al.,
``From AI for Science to Agentic Science:
A Survey on Autonomous Scientific Discovery,''
arXiv:2508.14111 (2025).

\bibitem[Douglas(2022)]{Douglas2022}
M.~R.~Douglas,
``Machine learning as a tool in theoretical science,''
Nature Rev.\ Phys.\ \textbf{4}, 145--146 (2022).

\bibitem[Brunton et al.(2016)]{Brunton2016}
S.~L.~Brunton, J.~L.~Proctor and J.~N.~Kutz,
``Discovering governing equations from data by sparse identification of
nonlinear dynamical systems,''
Proc.\ Natl.\ Acad.\ Sci.\ \textbf{113}, 3932--3937 (2016).

\bibitem[Schmidt and Lipson(2009)]{Schmidt2009}
M.~Schmidt and H.~Lipson,
``Distilling free-form natural laws from experimental data,''
Science \textbf{324}, 81--85 (2009).

\bibitem[Udrescu and Tegmark(2020)]{Udrescu2020}
S.-M.~Udrescu and M.~Tegmark,
``AI Feynman: A physics-inspired method for symbolic regression,''
Sci.\ Adv.\ \textbf{6}, eaay2631 (2020).

\bibitem[Chen et al.(2018)]{Chen2018}
R.~T.~Q.~Chen, Y.~Rubanova, J.~Bettencourt and D.~Duvenaud,
``Neural ordinary differential equations,''
Advances in Neural Information Processing Systems \textbf{31} (2018).

\bibitem[Kovachki et al.(2023)]{Kovachki2023}
N.~Kovachki et al.,
``Neural operator: Learning maps between function spaces,''
Acta Numerica \textbf{31}, 1--145 (2023).

\bibitem[Li et al.(2021)]{Li2021}
Z.~Li et al.,
``Fourier neural operator for parametric partial differential equations,''
arXiv:2010.08895 (2021).

\bibitem[Lu et al.(2021)]{Lu2021}
L.~Lu, P.~Jin and G.~E.~Karniadakis,
``Learning nonlinear operators via DeepONet based on the universal
approximation theorem of operators,''
Nature Machine Intelligence \textbf{3}, 218--229 (2021).

\bibitem[Veli\v{c}kovi\'c et al.(2018)]{Velickovic2018}
P.~Veli\v{c}kovi\'c et al.,
``Graph attention networks,''
International Conference on Learning Representations (2018).

\bibitem[Bronstein et al.(2021)]{Bronstein2021}
M.~M.~Bronstein, J.~Bruna, Y.~Cohen and P.~Veli\v{c}kovi\'c,
``Geometric deep learning: Grids, groups, graphs, geodesics, and gauges,''
arXiv:2104.13478 (2021).

\bibitem[No\'e et al.(2020)]{Noe2020}
F.~No\'e, S.~Olsson, J.~K\"ohler and H.~Wu,
``Boltzmann generators: Sampling equilibrium states of many-body systems
with deep learning,''
Science \textbf{365}, eaaw1147 (2019).

\bibitem[Carleo et al.(2019b)]{Carleo2019b}
G.~Carleo et al.,
``Machine learning and the physical sciences,''
Rev.\ Mod.\ Phys.\ \textbf{91}, 045002 (2019).

\end{thebibliography}

\printindex

\end{document}
